% !TeX root = ../papers/nc-p3.tex
\section{二辺だけの二元商から対称性を仮定せず四窓を復元する}
\label{sec:binary-two-edge-rigidity}
\begin{theorem}\label{thm:binary-two-edge-rigidity}
標準有限Frobenius四窓の連続二辺に、両端の文字をそれぞれ三次元に割る二元形式の商があるとする。
中央の四文字は両方の商で三次二元形式の双対となる。両辺には純関係がないとする。
このとき中央に現れる二つのねじれ三次曲線は自動的に一致し、窓全体が
定理\ref{thm:binary-apolar-rank-test}の階数2の両対称非正則窓となる。
最初から二辺の二元座標の整合性、四周期の印、あるいは交換対称性を仮定する必要はない。
特に全辺に有限点グラフを持つFrobenius窓では、この隣接二辺の検査だけで、
交互の二元形式による指数成長商の存在が必要十分に判定できる（始点は動かす）。
そのような商を持つ窓は全て両対称で反対階数8、二元座標の一周の移動は恒等である。
\end{theorem}
\begin{proof}[中央の三次曲線が異なる場合の圧縮交叉]
各辺の六関係は端の一文字を割っても単射で、双対の積は端の三平面では二元乗法である。
三次交叉を両端で同時に圧縮しても単射となる。
実際三重形式を二つの三平面と中央文字で知れば、第一辺の関係の単射性から第一端全体へ、
次に第二辺の関係の単射性から最終端全体へ一意に戻る。
従って圧縮された二辺の三次交叉 $S$ は少なくとも四次元である。

中央の二曲線 $C,C'\subset\mathbb P^3$ が相異なると仮定し、交叉長を $m$ とする。
$m\ge6$ なら $C'$ の三つの二次式は $C$ 上で高々一次元の六次式空間に制限される。
従って二つの独立な二次式が両曲線を含む。非退化整曲線を含む二次式は既約なので、
次数4の完全交叉が次数6の和を含むことになり矛盾。よって $m\le5$。
$C$ の座標を $x$、$C'$ の座標を $y=[s:t]$ とする。$C'$ の標準座標で $C$ を
独立な三次形式 $a_0(x),\ldots,a_3(x)$ によって表示する。
$S^*$ は $P_5(x)\otimes P_2(y)$ を
\[
 F_i(x,y)=a_i(x)t-a_{i+1}(x)s\qquad (0\le i\le2)
\]
の $P_2(x)\otimes P_1(y)$ 倍で割ったものになる。これは左辺でまず二元乗法を取り、
右辺の乗法の核を対角式の倍で表示したものである。関係を入れる前後がともに18次元なので、
$\dim S$ は $\sum b_i F_i=0$ を満たす双次数 $(2,1)$ の三係数の空間の次元でもある。

$F_i=u_i s+v_i t$ とし、$C'$ の両座標端点を交叉から外して $u,v$ をともに基点自由な三次式の組にする。
$b_i=A_i s+B_i t$ と置けば $A\cdot u=B\cdot v=0$ で、残る条件は
$A\cdot v+B\cdot u=0$。最初の二つの解空間は各三次元である。
実際、各基点自由空間には互いに素な二つの三次式があり、その二次式倍だけで $P_5$ を張る。
$u\times v$ の成分は $C'$ の三つの二次式の制限であり、共通因子 $d$ はちょうど交叉を切り、次数 $m$。
両方に直交する二次式列は原始列 $(u\times v)/d$ の多項式倍だから、その次元は
$\max(3-n,0)$、ただし $n=6-m$。$A\cdot v,B\cdot u$ はともに $dP_{5-m}$ に入る。
これは $u$ のある成分が各点の局所環で単元で、$d$ を法として $v$ は $u$ の倍数であることから分かる。
従って $m\le2$ なら両写像は単射で最後の一致の空間は高々三次元。
$m\ge3$ なら各像は $n$ 次元の $dP_{5-m}$ に全射し、$\dim S=6-n=m$ となる。

よって残る $m=4,5$ で $S$ は交叉スキーム $D=C\cap C'$ 上の三文字評価の双対像に等しい。
実際この像は両辺のHankel条件を満たし、$P_2(x)P_3(x)=P_5(x)$ は $D$ に全射、
残る二次文字にも $D$ 全体で単元なものを取れるから次元は $m$ である。
\end{proof}
\begin{proof}[非被約交叉を含む二つの追加条件]
第一端の第四文字での積を $L:P_3(x)\to P_5(x)$ と書く。
$D$ 上の一点評価をこの文字まで戻し、第二端を圧縮したまま取ると、中央と第二端には
$L(-)(x)\otimes\operatorname{ev}_{y}^{(2)}$ という値が残る。
これが右辺のHankel関係に入る条件をその零化空間に評価する。
右辺の乗法の核は二組の二元変数の対角式と双次数 $(2,1)$ の全空間の積なので、
得られる制限写像の双対は $D$ 上
\[
                b(x)\otimes\operatorname{ev}^{(1)}_{y}
\]
という評価列を与える。ここで $b$ は中央の二次式に、$y$ で消える一次式を掛けてから
$L(-)(x)$ を評価した三成分である。各台点で $b$ は非零である。
消えれば $L(-)(x)$ が中央の $y$ における評価の倍数となり、
$C$ と $C'$ の同じ点であるため第一辺の関係に純テンソルを与えてしまう。
線束を $D$ 上で自明化し、$b$ の一般の一次結合を取れば全台で単元になる。
従ってこの追加条件の階数は $P_1(y)$ の $D$ への制限階数2以上である。
$D$ は $C'$ 上の長さ4または5の部分スキームだから、全てが一点に支持されても階数2である。
これで持ち上がる三次交叉の次元は $m-2\le3$ となり矛盾する。
\end{proof}
\begin{proof}[一致した印からの全窓の復元]
二つの曲線が一致したので共通の二元座標を選び、両端での第四文字の積を
$L,\widetilde L:P_3\to P_5$ と書く。両端を圧縮した交叉は今度は全指数和だけに依存する $P_7^*$ である。
各持上げ $L$ の三つの差分 $L(x^{j+1})-xL(x^j)$ を六次に斉次化し、
両持上げの合計六元の張る空間を $D_6\subset P_6$ とする。
片端を第四文字に戻した値が他端でHankelである条件から、持ち上がる $D\in P_7^*$ は
$P_1D_6$ 全体を消す。従って $\dim(P_1D_6)\le4$ である。
各持上げの差分は基点自由である。共通零点では $L(-)$ の評価が三次評価の倍数となり、
その点での $P_5$ の評価は元の関係の中で純テンソルになるからである。
$\dim D_6\ge3$ なら一次式倍で少なくとも一次元増えるので、可能なのは次元3から4への極小増加だけ。
このとき $D_6=hP_2$ となる（命題\ref{prop:binary-quotient-rank-two-curve}の
$U+xU$ の証明）。共通因子があるため不可能である。
よって $D_6$ は互いに素な二つの六次式の鉛筆で、両辺の三差分はそれぞれこの鉛筆を張る。

二元多項式環をこの正則列で割るとソクル次数10のArtin完全交叉を得る。
その最上次数の汎関数 $T$ は $D_6P_4$ を消す唯一の射影汎関数で、
五次までは関係がないので $G=H_{5,5}(T)$ は可逆である。
差分を $P_4$ 倍して $T$ に入れれば $GL,G\widetilde L$ はそれぞれHankelとなる。
$P_1D_6$ の零化空間は四次元で $T\mathbin{\lrcorner}P_3$ に一致するので、全ての
$T\mathbin{\lrcorner}q'$ が実際に三次交叉に持ち上がる。
両端を第四文字にして左右から計算すると、一方は $G(Lq,\widetilde Lq')$、
他方は $G(Lq',\widetilde Lq)$ となる。実際片端だけ第四文字にすると
$T(pq' Lq)=T(pq Lq')$ のHankel等式がその側の関係係数を一意に決める。
従って混合行列は対称である。前定理の差分計算から、独立な二行を持つ二つの八次パラメータは
極性空間を法として比例する。係数は非零なので、第四文字を選び直せば両持上げは一致する。
これで三次交叉が前節のGram窓の四つの縮約で張られた。
元の窓も第四文字との対合が完全なので、第四位置の基底変更だけでGram窓に一致する。
全四辺の一意の二元印は同じ二元座標で実現され、一周の移動も恒等となる。
圧縮交叉の有理数標本と長さ5の一点における検算は
\path{computations/verify_binary_two_edge_rigidity.py}で再現できる。
\end{proof}
